% -*- coding: utf-8 -*-
% Issue #859 / #827 统一视觉对比 MWE
% 用三个版本的 xeCJK 分别编译此文件，对比视觉效果

\documentclass{article}
\usepackage{xeCJK}
\usepackage[paperwidth=21cm,paperheight=29cm,margin=1.2cm]{geometry}
\usepackage{tabularray}
\usepackage{xcolor}
\usepackage{enumitem}

\setmonofont{DejaVu Sans Mono}[Scale=0.75]
\parindent=0pt
\pagestyle{empty}
\hbadness=10000
\hfuzz=200pt
\setlength{\fboxsep}{1pt}

% 宽度对比工具
\newcommand{\wtest}[1]{%
  \begingroup
  \setbox0=\hbox{#1}%
  \setbox2=\vbox{\hsize=50cm\parfillskip=0pt\leftskip=0pt\rightskip=0pt #1\par}%
  \setbox4=\vbox{\unvbox2\unskip\unpenalty\global\setbox6=\lastbox}%
  \setbox8=\hbox{\unhbox6}%
  \makebox[5.5em][l]{\small hbox:\the\wd0}%
  \makebox[5.5em][l]{\small par:\the\wd8}%
  \ifdim\wd0=\wd8
    \textcolor{green!50!black}{\small\bfseries PASS}%
  \else
    \textcolor{red}{\small\bfseries FAIL}%
    {\small\enspace$\Delta$=\the\dimexpr\wd0-\wd8\relax}%
  \fi
  \endgroup
}

\begin{document}

% ====================================================================
{\large\bfseries\sffamily 1. 原始 Bug 复现}
\smallskip\hrule\smallskip

{\small\sffamily Issue \#859: tabularray 中右对齐+标点结尾单元格左侧多余间距}

\medskip
\begin{tblr}{columns={colsep=0pt}, columns={halign=r}, hlines, vlines}
  test  & 123456       \\
  test: & 左侧英文标点 \\
\end{tblr}
\quad
\begin{tblr}{columns={colsep=0pt}, columns={halign=r}, hlines, vlines}
  test    & 123456             \\
  test：  & 左侧中文标点全角式 \\
\end{tblr}
\quad
\xeCJKsetup{PunctStyle=plain}%
\begin{tblr}{columns={colsep=0pt}, columns={halign=r}, hlines, vlines}
  test   & 123456           \\
  test： & 左侧中文标点原样 \\
\end{tblr}

{\small\sffamily Issue \#827（comment）: tabular 中右对齐+标点结尾单元格左侧多余间距}

\begin{tabular}{|@{}r@{}|@{}r@{}|}
test  & 123456       \\
test: & 左侧英文标点 \\
\end{tabular}
\quad
\begin{tabular}{|@{}r@{}|@{}r@{}|}
test    & 123456             \\
test：  & 左侧中文标点全角式 \\
\end{tabular}
\quad
\xeCJKsetup{PunctStyle=plain}
\begin{tabular}{|@{}r@{}|@{}r@{}|}
test   & 123456           \\
test： & 左侧中文标点原样 \\
\end{tabular}

\bigskip

\xeCJKsetup{PunctStyle=kaiming}%

% ====================================================================
{\large\bfseries\sffamily 2. 宽度一致性验证}
\smallskip\hrule\smallskip

\begin{minipage}[t]{0.48\linewidth}
{\small\sffamily 以右标点结尾（\#859 核心场景）：}

\smallskip
\makebox[8em][l]{测试。}\wtest{测试。}

\makebox[8em][l]{测试，}\wtest{测试，}

\makebox[8em][l]{测试：}\wtest{测试：}

\makebox[8em][l]{测试"}\wtest{测试"}

\makebox[8em][l]{测试》}\wtest{测试》}
\end{minipage}%
\hfill
\begin{minipage}[t]{0.48\linewidth}
{\small\sffamily 对照组（应始终 PASS）：}

\smallskip
\makebox[8em][l]{测试内容}\wtest{测试内容}

\makebox[8em][l]{测试hello}\wtest{测试hello}

\makebox[8em][l]{测试123}\wtest{测试123}

\makebox[8em][l]{测试，内容}\wtest{测试，内容}

\makebox[8em][l]{(空)}\wtest{测试}
\end{minipage}

\bigskip

% ====================================================================
{\large\bfseries\sffamily 3. 方案 1 影响范围：段中 Boundary 转换（heading 类）}
\smallskip\hrule\smallskip

{\small 方案 1 在标点后接非 CJK 字符处多一个 \texttt{\textbackslash penalty 0}。
方案 2 和修复前在此处节点相同。以下验证视觉效果：}

\medskip
\begin{minipage}[t]{0.48\linewidth}
\fbox{\parbox{16em}{%
{\bfseries\sffamily 第一章：概述}

本章概述全文内容。
}}
\end{minipage}%
\hfill
\begin{minipage}[t]{0.48\linewidth}
\fbox{\parbox{16em}{%
{\bfseries\sffamily 问题：关于 xeCJK}

标题中混合使用标点和字母。
}}
\end{minipage}

\medskip
\fbox{\parbox{0.96\linewidth}{%
正文段落：hello world 继续中文。第一条：100 个单位。括号（parenthesis）混排。%
}}

\bigskip

% ====================================================================
{\large\bfseries\sffamily 4. 方案 2 影响范围：段末标点（basic/beamer 类）}
\smallskip\hrule\smallskip

{\small 方案 2 在段末多一个 \texttt{glue + penalty 0}。
方案 1 在段末也保留了 glue 和 penalty 0（方式不同但结果相同）。
以下验证视觉效果：}

\medskip
\begin{minipage}[t]{0.48\linewidth}
\fbox{\parbox{14em}{%
天地玄黄，宇宙洪荒。日月盈昃，辰宿列张。
}}
\end{minipage}%
\hfill
\begin{minipage}[t]{0.48\linewidth}
\fbox{\parbox{14em}{%
设 $f(x) = x^2$，则 $f(0) = 0$。

以上公式成立。
}}
\end{minipage}

\medskip
\fbox{\parbox{0.96\linewidth}{%
这是一段以右书名号结尾的文字，引用《西游记》%
}}

\fbox{\parbox{0.96\linewidth}{%
Beamer 幻灯片标题中使用 ctexbeamer 排版中文报告。
}}

\bigskip

% ====================================================================
{\large\bfseries\sffamily 5. 极端 Bad Case：窄行宽断行压力}
\smallskip\hrule\smallskip

{\small 方案 1 的 \texttt{\textbackslash penalty 0} 理论上可能在极端窄行宽下
影响 \TeX{} 的最优断行。以下用 5em 行宽测试：}

\medskip
\begin{minipage}[t]{0.22\linewidth}
{\small\sffamily 5em 行宽：}

\fbox{\parbox{5em}{%
这是测试：hello继续中文。
}}
\end{minipage}%
\hfill
\begin{minipage}[t]{0.22\linewidth}
{\small\sffamily 4em 行宽：}

\fbox{\parbox{4em}{%
问题：AB和CD结论。
}}
\end{minipage}%
\hfill
\begin{minipage}[t]{0.22\linewidth}
{\small\sffamily 6em 行宽：}

\fbox{\parbox{6em}{%
列表：第一，第二，第三。
}}
\end{minipage}%
\hfill
\begin{minipage}[t]{0.22\linewidth}
{\small\sffamily 5em 行宽：}

\fbox{\parbox{5em}{%
"引号"后text续排。
}}
\end{minipage}

\medskip
{\small\sffamily 极端窄行宽 (3em)：}

\smallskip
\fbox{\parbox{3em}{测试：A。}}\enspace
\fbox{\parbox{3em}{问题，B。}}\enspace
\fbox{\parbox{3em}{结论：C！}}\enspace
\fbox{\parbox{3em}{正常text}}

\end{document}
